% =====================================================================
% Appendix E --- Pricing extensions: input demands, accounting and robustness
% ---------------------------------------------------------------------
% Job: the input-demand equations and profit accounting for Section 6,
%      with numerical checks and the all-sector productivity counterfactual.
% Code: OpenCGE@src/closures/pricing/pricing.jl, OpenCGE@src/core/equilibrium.jl
%      (_mobile_market_demand, _cost_minimizing_labor), OpenCGE@src/core/
%      diagnostics.jl (gdp_components, retained_profit); the gate lives in
%      experiments/run.jl (evaluate_gates). Records: ADR-0029, ADR-0030, ADR-0031.
% =====================================================================
\section{Pricing extensions: input demands, accounting and robustness}\label{app:pricing-extensions}

This appendix gives the input demands and profit-accounting identity for the
pricing extensions in Section~\ref{sec:beyond-demand-only}, followed by the
results when productivity effects extend to all sectors.

\subsection{Input demands}\label{app:pricing-hooks}
Sector $i$ has the CES unit cost
\begin{equation}\label{eq:app-unit-cost}
  c_i \;=\; \Bigl[\,(\mathcal{A}_i^{H})^{\varepsilon-1}\bigl(\alpha_i\,w_i^{1-\varepsilon}
     + (1-\alpha_i)\,P_i^{1-\varepsilon}\bigr)\Bigr]^{\frac{1}{1-\varepsilon}}.
\end{equation}
A firm producing $q$ incurs cost $q\,c_i(p,w_i;\mathcal{A}_i^{H}(y_i))$, taking aggregate sector output $y_i$ as given. Its marginal cost is therefore $c_i$, and the input-price derivatives below hold effective productivity fixed.
At given output and effective productivity, Shephard's lemma applied to \eqref{eq:app-unit-cost} gives labor demand and, through the domestic fraction $a_i/(1-\alpha_i)$ of the composite (Section~\ref{subsec:optimization}), domestic intermediate demand:

\begin{align}
  L_i \;&=\; c_i^{\varepsilon}\,(\mathcal{A}_i^{H})^{\varepsilon-1}\,\alpha_i\,y_i\,w_i^{-\varepsilon},
  \label{eq:app-labor}\\
  x_{ui} \;&=\; p_u^{-\theta}\,\Omega_{iu}\;c_i^{\varepsilon}\,
     (\mathcal{A}_i^{H})^{\varepsilon-1}\,P_i^{\theta-\varepsilon}\,a_i\,y_i ,
  \label{eq:app-intermediate}
\end{align}

Firms pay wages $w_i$ and input prices $p_u$; the unit-cost term in these
demands is marginal cost $c_i=p_i/\mu_i$. The common factor
$k_i=c_i^{\varepsilon}(\mathcal{A}_i^{H})^{\varepsilon-1}P_i^{1-\varepsilon}$ also values
imported intermediates $M^{\mathrm{int}}$ and product taxes on intermediate use
$T^{\mathrm{int}}$ in the accounts. The price therefore enters the first-order
conditions only through marginal cost: writing $p_i$ for $c_i$ in any demand term
inflates it by $\mu_i^{\varepsilon}$ while leaving the price response untouched.

\paragraph{Checks.} Every extension is reached by continuation from the
demand-only model and re-solved from four starting points; a cell is accepted only if the four solutions agree to within $10^{-6}$ and the residual Jacobian has full rank. At
$\delta=\kappa=0$ every priced cell reproduces its cell of
Table~\ref{tab:results-matrix}, which is what lets the extension's incremental effects be calculated by
subtracting the corresponding zero-extension row.

\subsection{Profit income and accounting}\label{app:pricing-accounting}
Distributed profits enter household income,
$E=(1-\tau)\sum_i w_iL_i+F+(1-s_\pi)\Pi$, where $s_\pi$ is the share of profits
retained before household saving, while retained profits enter the external account:
\begin{equation}\label{eq:app-account}
  sE+s_\pi\Pi+T^{\mathrm{int}}+M-(I+X)
  = F+B_{\mathrm{gov}}.
\end{equation}

Income-side GDP is $\sum_i w_iL_i+\Pi$. In the three closures whose numeraire
is the CPI --- no reallocation, standard CGE and elastic supply --- F2 and
F3 remain indistinguishable in every row, so the transfer neutrality behind
\eqref{eq:finstrument} survives the pricing extension.


\subsection{Productivity effects in all sectors}\label{app:profiles}
The main text applies the utilization elasticity to the goods sectors only. In
the classification of the source table (NACE, CPA sections) these are sections
A--F: agriculture, forestry and fishing; mining and quarrying; manufacturing;
electricity, gas, steam and air conditioning; water supply, sewerage and waste;
and construction. Services (sections G onward) keep $\delta = 0$ in the
main-text runs.

Table~\ref{tab:pricing-profiles} reports the all-sector counterfactual
introduced in Section~\ref{subsec:pricing-capacity}, extending the productivity
effects to services. The Leontief closure shows the same financing-dependent
reversal of the real-wage response ($\Delta\omega$) as the employment-adjustment closure. The
profile is therefore a statement about where returns are assumed to live, not
about the size of the coefficient.

\clearpage
% =====================================================================
% Table E --- pricing-door profiles not carried in the main text (GENERATED)
% ---------------------------------------------------------------------
% Do NOT edit by hand. Regenerate with:  make_pricing_door.py
% OpenCGE tag v0.5.0 (ADR-0044); source run ids are listed below.
% Sources: pricing_capacity_rows-<row>-<fin>-<dp05|dm05>,
% pricing_capacity-GAMMA-<fin>-<dp05|dm05>,
% pricing_mobile-ALPHA-<fin>-<dp05|dm05>.
% One panel: the all-sector capacity counterfactual; the kappa = 0.25
% distributed-profit rows were dropped from the paper on 2026-10-24
% (they repeated the invariance Section 6.2 already states).
% =====================================================================
\begin{table}[htbp]
\centering
\scriptsize
\setlength{\tabcolsep}{1.6pt}
\renewcommand{\arraystretch}{1.15}
\resizebox{\textwidth}{!}{%
\begin{tabular}{@{}l@{\hspace{6pt}}rrrrr@{\hspace{6pt}}rrrrr@{\hspace{6pt}}rrrrr@{}}
\toprule
 & \multicolumn{5}{c}{\textbf{F1}} & \multicolumn{5}{c}{\textbf{F2}} & \multicolumn{5}{c}{\textbf{F3}} \\
\cmidrule(lr){2-6}\cmidrule(lr){7-11}\cmidrule(lr){12-16}
\multicolumn{1}{c}{Labor closure and $\delta$} & \multicolumn{1}{c}{$\Delta\omega$} & \multicolumn{1}{c}{\shortstack{$\max$\\$|\Delta\hat p|$}} & \multicolumn{1}{c}{$\Delta L$} & \multicolumn{1}{c}{$\Delta W_D$} & \multicolumn{1}{c}{$\frac{\Delta\Pi}{Y}$} & \multicolumn{1}{c}{$\Delta\omega$} & \multicolumn{1}{c}{\shortstack{$\max$\\$|\Delta\hat p|$}} & \multicolumn{1}{c}{$\Delta L$} & \multicolumn{1}{c}{$\Delta W_D$} & \multicolumn{1}{c}{$\frac{\Delta\Pi}{Y}$} & \multicolumn{1}{c}{$\Delta\omega$} & \multicolumn{1}{c}{\shortstack{$\max$\\$|\Delta\hat p|$}} & \multicolumn{1}{c}{$\Delta L$} & \multicolumn{1}{c}{$\Delta W_D$} & \multicolumn{1}{c}{$\frac{\Delta\Pi}{Y}$} \\
\midrule
\multicolumn{16}{@{}l}{\textit{Capacity effects on all sectors: No-reallocation}} \\
\quad $\delta = +0.5$ & $+0.73$ & $26.82$ & $0.00$ & $-0.01$ & $0.00$ & $+0.85$ & $35.47$ & $0.00$ & $-0.03$ & $0.00$ & $+0.85$ & $35.47$ & $0.00$ & $-0.03$ & $0.00$ \\
\quad $\delta = -0.5$ & $+0.46$ & $16.39$ & $0.00$ & $-0.03$ & $0.00$ & $+0.50$ & $19.52$ & $0.00$ & $-0.04$ & $0.00$ & $+0.50$ & $19.52$ & $0.00$ & $-0.04$ & $0.00$ \\
\midrule
\multicolumn{16}{@{}l}{\textit{Capacity effects on all sectors: Standard CGE}} \\
\quad $\delta = +0.5$ & $+0.12$ & $4.06$ & $0.00$ & $+0.07$ & $0.00$ & $+0.13$ & $4.27$ & $0.00$ & $+0.07$ & $0.00$ & $+0.13$ & $4.27$ & $0.00$ & $+0.07$ & $0.00$ \\
\quad $\delta = -0.5$ & $-0.18$ & $4.56$ & $0.00$ & $+0.04$ & $0.00$ & $-0.18$ & $4.47$ & $0.00$ & $+0.04$ & $0.00$ & $-0.18$ & $4.47$ & $0.00$ & $+0.04$ & $0.00$ \\
\midrule
\multicolumn{16}{@{}l}{\textit{Capacity effects on all sectors: Elastic labor supply ($\eta_s = 0.5$)}} \\
\quad $\delta = +0.5$ & $+0.33$ & $14.11$ & $+0.15$ & $+0.21$ & $0.00$ & $+0.36$ & $16.21$ & $+0.16$ & $+0.22$ & $0.00$ & $+0.36$ & $16.21$ & $+0.16$ & $+0.22$ & $0.00$ \\
\quad $\delta = -0.5$ & $+0.17$ & $4.31$ & $+0.07$ & $+0.26$ & $0.00$ & $+0.18$ & $4.54$ & $+0.08$ & $+0.27$ & $0.00$ & $+0.18$ & $4.54$ & $+0.08$ & $+0.27$ & $0.00$ \\
\midrule
\multicolumn{16}{@{}l}{\textit{Capacity effects on all sectors: Employment adjustment}} \\
\quad $\delta = +0.5$ & $+0.10$ & $4.06$ & $+0.04$ & $+0.12$ & $0.00$ & $-0.03$ & $4.22$ & $+0.35$ & $+0.48$ & $0.00$ & $-0.99$ & $3.90$ & $+2.61$ & $+3.10$ & $0.00$ \\
\quad $\delta = -0.5$ & $-0.34$ & $4.60$ & $-0.15$ & $-0.42$ & $0.00$ & $-0.19$ & $4.47$ & $-0.01$ & $+0.02$ & $0.00$ & $+0.94$ & $4.03$ & $+1.05$ & $+3.32$ & $0.00$ \\
\midrule
\multicolumn{16}{@{}l}{\textit{Capacity effects on all sectors: Leontief closure}} \\
\quad $\delta = +0.5$ & $+0.12$ & $4.52$ & $+0.02$ & $+0.10$ & $0.00$ & $-0.02$ & $4.54$ & $+0.31$ & $+0.44$ & $0.00$ & $-1.04$ & $4.21$ & $+2.45$ & $+2.92$ & $0.00$ \\
\quad $\delta = -0.5$ & $-0.28$ & $4.12$ & $-0.12$ & $-0.35$ & $0.00$ & $-0.14$ & $4.15$ & $+0.03$ & $+0.10$ & $0.00$ & $+0.90$ & $3.76$ & $+1.22$ & $+3.46$ & $0.00$ \\
\bottomrule
\end{tabular}}
\caption{Kaldor--Verdoorn effects extended to every sector ---
the all-sector counterfactual of Table~\ref{tab:pricing-kaldor} --- on all five
labor closures. Conventions, including the real-wage and dispersion conventions
of Table~\ref{tab:pricing-kaldor}, as in Table~\ref{tab:pricing-markup}.}
\label{tab:pricing-profiles}
\end{table}

